\section{Data and Empirical Design}
\label{sec:empirical-design}

\subsection{Data and timing}
\label{sec:panel-data}

The design evaluates the financial relevance of peer relations measured from
an earlier information window.
We constructed the barycentric field and every comparator field from
information ending in 2022 and froze each matrix.
We then evaluated conditional return dependence on
\ppvalue{post-pooled-2023-2026-w-flat-trading-days} common trading days from
\ppvalue{post-evaluation-date-start} through
\ppvalue{post-evaluation-date-end}.
No evaluation return enters the text geometry or any comparator matrix.
Return data serve two roles.
The reported post-corpus QMLE estimates are fitted on 2023--2026 returns, while
the held-out density exercise fits every parameter on 2018--2022 returns and
evaluates it on 2023--2026.

We assembled firm text from Nasdaq's public per-symbol archive of syndicated
news \citep{nasdaq_news_headlines_2026}, whose ticker-indexed retrieval and
dated article bodies support one reproducible collection rule across the
universe.
Ticker assignment follows the archive and is not manually validated at the
article level.
Daily simple returns are Yahoo Finance adjusted-close ratios minus one,
retrieved through a version-fixed \texttt{yfinance} client
\citep{yahoo_finance_historical_2026,aroussi_yfinance_2026}.
The primary specification subtracts the daily risk-free rate and aligns the
excess-return panel with the daily Fama--French five factors and momentum.
We keep only dates on which every firm has a return; missing returns are not filled.

The universe is a prespecified Nasdaq-100-based frame admitted by an annual
coverage screen requiring at least \ppvalue{annual-raw-article-floor} raw
articles per firm in each year from 2018 through 2022.
Intersecting firms that have a frozen embedding cloud with firms that have a
compatible return history yields \ppvalue{n-tickers-priced} continuously and
heavily covered large-cap firms, a survivor-conditioned cross-section.
The annual floor preserves coverage in every construction year; both that
floor and the balanced cloud size below were fixed before the return
comparisons.
The return window ends on 6 April 2026, the final common date before HOLX's
take-private,
so the 2026 fold covers
\ppvalue{post-eval2026-w-flat-trading-days} dates while the 2023--2025 folds
cover complete years.
\ppAppendixRef{sec:app-sample} reports sector composition and bridges the
estimates to an earlier 52-firm specification.

Each firm's information environment is a balanced cloud of
\ppvalue{paper5-balanced-cloud-size} articles, selected from the 2018--2022
window by a deterministic URL-hash order.
The common cloud size puts every empirical law on the support required by the
balanced assignment in \eqref{eq:w2-assignment}; it makes firms comparable and,
by design, removes raw news volume from the geometry.
It also fixes the assignment problem's computational support across firms.
We encoded each article with the full-width Qwen3-Embedding-8B model, which is
trained to place semantically related texts near one another
\citep{qwen3_embedding_2025}, and row-normalized its 4,096-coordinate vector.
The encoder was released after the 2018--2022 article window, so the primary
representation is a retrospective measurement of earlier news.
The earlier-cutoff comparison in \Cref{sec:results-representation}
examines dependence on this retrospective measurement.
Each article has equal mass in the firm's empirical distribution.
The representation exercises in \Cref{sec:results-representation} vary output
width, model capacity, model family, and reported information vintage
(\ppAppendixRef{sec:p5-representation-ablation}).

\subsection{Estimation, inference, and out-of-period evaluation}
\label{sec:estimation}

We estimate conditional return dependence by quasi-maximum likelihood
(QMLE).
The primary specification is
\eqref{eq:factor-spatial-lag} with the barycentric field $W^\flat$, firm-specific
intercepts, and loadings on the market, size, value,
profitability, investment, and momentum factors.
At every candidate $\rho$, least squares profiles all firm intercepts and
loadings jointly with $\sigma^2$, and the concentrated objective retains the
spatial Jacobian $T\log|I-\rho W|$.
The search covers the stable region $|\rho|<1$.
The factor-only benchmark sets $\rho=0$ with the same intercepts, loadings,
dates, and innovation specification, so comparisons with it isolate the value
of the field.
Following \citet{ge_news-implied_2023}, a control ladder repeats the primary
specification under FF3, FF3 plus momentum, and FF5 on the same dates; FF5 plus
momentum remains primary throughout.
A robustness specification profiles firm-specific innovation variances
$D=\operatorname{diag}(\sigma_1^2,\ldots,\sigma_N^2)$, because ignored
heteroskedasticity can move a spatial point estimate as well as its standard
error \citep{kelejian_prucha_2010}.
Maximized quasi-log-likelihood (QLL) describes fit on the estimation panel.
The appendix reports the associated working AIC, Hannan--Quinn (HQC), and BIC
criteria, their definitions, and parameter counts
(\ppAppendixRef{sec:app-information-criteria}).

Two bootstrap procedures supply inference. A restricted-null stationary
bootstrap tests $H_0:\rho=0$ through the quasi-likelihood-ratio (QLR) statistic
by independently resampling standardized firm residuals on stationary-block
paths, restoring the fitted variances and factor means before refitting both
models. Interval estimates and every held-out comparison use 2,000 paired
date-block refits on one common resampling schedule, which refits the
calibration model, the evaluation coefficients, and the scored parameters on
each draw; all models and both transport orders share every draw. Both
procedures condition on the corpus, embeddings, and constructed fields, so they
propagate return sampling and nuisance-parameter uncertainty. Implementation
details are given in \ppAppendixRef{sec:app-inference}.

The main financial evaluation asks whether the field improves the joint
distribution assigned to later returns.
All parameters are fitted on 2018--2022 returns and held fixed while 2023--2026
returns are scored by the mean Gaussian quasi-log score per firm-day.
Higher scores mean that the model assigns greater density to the realized
cross-section of returns.
Evaluation-date factors are observed when scoring, so the comparison concerns
conditional joint density given realized factor returns.
Paired intervals resample calibration and evaluation dates on separate
stationary-block schedules shared by all models.

\subsection{Transport order and the comparator fields}
\label{sec:two-field-estimation}

The transport-order comparison repeats the target-anchored construction under
linear ($p=1$) and quadratic ($p=2$) displacement cost, holding fixed the firm
clouds, ground space, leave-one-out candidates, and simplex.
For $W^\flat_1$, each row minimizes the sum of Euclidean reconstruction norms
after order-one optimal assignments; for $W^\flat$, the sum of squared norms
after order-two assignments.
Tables and the comparison below label the two orders W1 and W2.
Changing the order changes how embedding-space displacement enters the
measurement of the field; the exposure-adjustment criterion that maps $\rho$
into $\lambda$ is the same for both, so the comparison measures sensitivity to
the transport geometry rather than the curvature of adjustment costs.
Before the final estimation we fixed an equivalence margin for the
W2-minus-W1 held-out score and a weak-separation rule that classifies the
fields as nearly collinear from the correlation of their induced peer-return
signals; when the separation rule binds, the design reports paired contrasts
rather than interpreting a joint coefficient allocation. Both thresholds and
their construction are stated in \ppAppendixRef{sec:app-inference}.
A Matryoshka
dimension ladder repeats both constructions at
$D\in\{4096,1024,256,64\}$ coordinates of the same model
\citep{kusupati_matryoshka_2022,qwen3_embedding_2025} to locate where any
difference between the orders enters (\ppAppendixRef{sec:p5-matryoshka}).

Each comparator field changes one element of the peer definition, holding the
firms, dates, factor controls, and likelihood fixed.

\begin{table}[H]
  \centering
  \caption{Comparator specifications and the element of the dependence
    definition each one changes. All fields use firm information dated through
    2022 and are held fixed throughout the 2023--2026 evaluation; covariance
  inputs are calibrated on 2018--2022 returns.}
  \label{tab:comparator-design}
  \small
  \begin{tabular}{p{0.3\linewidth}p{0.62\linewidth}}
    \toprule
    Specification & Relative to the barycentric field \\
    \midrule
    Equal active support & Same peers; equal instead of fitted weights \\
    RBF--Wasserstein (median or validation-tuned bandwidth) & Same
    distributions and pairwise distances; separate distance-decay weight per
    pair instead of joint reconstruction \\
    Centroid reconstruction & Same simplex reconstruction; each firm's cloud
    collapsed to its mean embedding \\
    Persistent news co-mentions & Explicit shared coverage instead of
    language-model geometry \\
    Industry/sector & Conventional classification \\
    Direct Ledoit--Wolf residual covariance & No peer field; constant-correlation
    shrinkage of the factor-residual covariance \\
    W1 barycentric & Linear instead of quadratic transport cost \\
    \bottomrule
  \end{tabular}
\end{table}

The equal-support, RBF, and co-mention fields are defined in
\Cref{sec:interaction-comparators}.
The centroid field solves the leave-one-out simplex problem on mean
embeddings, and the tuned RBF selects its bandwidth multiplier from
$\{0.25,0.5,1,2,4\}$ using only a 2018--2020 fit and 2021--2022 validation
split before refitting through 2022.
Because these fields are non-nested, their held-out gains are compared through
paired stationary-block intervals rather than QLR statistics.

As a structure-agnostic covariance benchmark, we estimate the Ledoit--Wolf
linear-shrinkage covariance of the factor residuals in the 2018--2022
calibration sample~\citep{ledoit_wolf_2004}. The estimator shrinks the sample
covariance toward a constant-correlation target that preserves each residual
variance and replaces pairwise correlations with their cross-sectional mean;
the shrinkage intensity is estimated from the calibration residuals. We hold
the factor fitted means and this covariance fixed when scoring the same
2023--2026 observations. This benchmark asks whether the information field
adds held-out density beyond a conventional covariance estimator without an
economic peer structure.

The relation between the barycentric and news-link fields is estimated jointly,
as the two-channel model of \Cref{thm:two-field-closure} prescribes, under the
same primary factor-conditioned return specification.
Scalar projection applies to \eqref{eq:two-field-sar} exactly as to the
one-field equilibrium, giving

\begin{equation}\label{eq:two-field-spatial-lag}
  (I-\rho^\flat W^\flat-\rho^{\mathrm{news}} W^{\mathrm{news}})r_t^e
  =\alpha+A f_t+\varepsilon_t .
\end{equation}

This comparison uses the firm-specific intercepts and factor loadings of the
primary model. Its matched-random overlap tests use the same frozen fields and
return panel.
The raw-return, common-intercept counterpart, which also supplies the annual
re-estimates, is reported in the appendix as a corroborating specification:
\[
  r_t=\alpha\1+\rho W r_t+\varepsilon_t,
  \qquad \varepsilon_t\mathrel{\overset{\mathrm{work}}{\sim}}
  \mathrm{iid}(0,\sigma^2I),
\]
extended to two fields as in \eqref{eq:two-field-spatial-lag}.
Both matrices enter unchanged: residualizing one on the other would create
negative entries and non-unit row sums and destroy the peer-average
interpretation of \Cref{def:interaction-matrix}, so the coefficient would no
longer map to an adjustment intensity.
Residualizing the induced signals $W^\flat r_t$ and $W^{\mathrm{news}}r_t$
instead would use the contemporaneous returns and so would not yield
predetermined fields for the spatial likelihood.
Because $I-\rho^\flat W^\flat-\rho^{\mathrm{news}}
W^{\mathrm{news}}=I-(\rho^\flat+\rho^{\mathrm{news}})W(\theta)$,
the concentrated objective is the one-field objective at the mixture, and the
restrictions $0\leq\theta\leq1$ and $0\leq\rho^\flat+\rho^{\mathrm{news}}<1$ impose
$\lambda^\flat,\lambda^{\mathrm{news}}\geq0$.
The single-field fits are the boundaries $\theta=1$ and $\theta=0$ of this
family, so the joint quasi-log-likelihood cannot fall below either.
The statistic
$\mathrm{QLR}^{\mathrm{news}}=2\{\sup_{\rho^\flat,\rho^{\mathrm{news}}\geq 0}
  \ell^Q(\rho^\flat,\rho^{\mathrm{news}})
-\sup_{\rho^{\mathrm{news}}=0}\ell^Q(\rho^\flat,\rho^{\mathrm{news}})\}$ tests
$\rho^{\mathrm{news}}=0$, $\mathrm{QLR}^\flat$ tests $\rho^\flat=0$
symmetrically, and each has its
reference distribution simulated under the restricted single-field fit because
the null lies on the boundary
(\ppAppendixRef{sec:p5-two-field-implementation}).
